\textbf{动态DP}

用于解决带修树上DP。核心思想：将状态转移写成广义矩阵乘法，用重链剖分+线段树维护重链上的矩阵连乘。

\textbf{状态抽象}

将节点 $u$ 的DP状态设为列向量 $F_u$。将所有轻儿子的贡献合并成转移矩阵 $G_u$，彻底剥离重儿子 $son_u$ 的转移：
\[ F_u = G_u \otimes F_{son_u} \]
此时重链顶端状态即为链上所有节点 $G$ 的连乘。

\textbf{线段树维护}

每条重链单独对应线段树区间，节点维护该点 $G_u$。
\begin{itemize}
    \item \textbf{合并方向}：若叶子按 dfn 序递增（即由浅到深），因 $F_u = G_u \otimes F_{son_u}$ 父亲在左，故合并顺序为 \texttt{左儿子 $\otimes$ 右儿子}。由于通常无交换律，方向绝不可错。
    \item \textbf{单点修改}：更新 $u$ 的 $G_u$。该修改改变了当前重链的总乘积，也就是改变了链顶 $top_u$ 传给 $fa_{top_u}$ 的轻儿子贡献。计算该贡献的变化量，更新 $fa_{top_u}$ 的 $G$ 并继续向上跳跃至根。时间 $O(\log^2 N)$。
\end{itemize}
